\chapter*{Introduzione}
\addcontentsline{toc}{chapter}{Introduzione}
\label{chap:introduzione}

Con il presente elaborato ci proponiamo di riassumere, analizzare in dettaglio e presentare criticamente i risultati emersi durante lo svolgimento delle tre esercitazioni pratiche del corso di \textbf{Sensori, Rivelatori e Dispositivi Elettronici}.

Per condurre una trattazione rigorosa ed efficace, adottiamo in tutto il documento un approccio metodologico strutturato in quattro fasi sequenziali:
\begin{enumerate}
    \item \textbf{Inquadramento teorico}: introduciamo anzitutto i fondamenti fisici e analitici strettamente necessari alla comprensione del fenomeno;
    \item \textbf{Scelte implementative}: discutiamo ed esplicitiamo le opzioni di modellazione numerica e di calcolo adottate;
    \item \textbf{Analisi dei risultati}: esaminiamo nel dettaglio le prestazioni e le risposte spettrali o statistiche ottenute;
    \item \textbf{Validazione e confronto}: mettiamo infine a confronto i dati simulati e sperimentali con i comportamenti teorici attesi per verificare la bontà dei modelli.
\end{enumerate}

Le tre esercitazioni si muovono lungo tre direttrici complementari dell'ingegneria dei sensori:
\begin{itemize}
    \item La prima affronta la \textbf{simulazione numerica 1D} di trasduttori piezoelettrici in ambiente MATLAB;
    \item La seconda estende l'analisi al dominio \textbf{FEM tridimensionale} mediante l'impiego del codice commerciale Ansys;
    \item La terza declina i concetti teorici nell'ambito dell'\textbf{analisi di caratteristiche biometriche} e dell'elaborazione di segnali ecografici del \textit{palm print} in MATLAB.
\end{itemize}

Sulla scorta di questa impostazione, procediamo ad addentrarci nei singoli argomenti trattati.
